\section{The quadratic disk coefficient}\label{sec:quadratic}

\begin{theorem}\label{thm:quadraticdisk}
For every Young diagram the two disk vertices, divided by their common
value at $y=0$, have the same coefficient of $y^2$ in $\QQ(p,c,x)$.
\end{theorem}

A reverse plane partition of size two has two distinct maximal boxes of
height one, a maximal box of height two, or a removable domino of height one.
We recover the first two terms by collecting the currents of each sign.
Interchanging two currents of opposite signs gives the domino term.

\subsection{Interchanging two currents}
Write $R=R_s$, and use the residues of \eqref{eq:onebody}. Thus $A_i^\alpha$ is the signed inner residue of $P_i(z)$ at
$sZ_\alpha$ and $B_i^\beta$ is the signed outer residue of $N_i(w)$ at
$Z_\beta$. Thus $\alpha$ is a downward step at or after $i$ and $\beta$
a horizontal step before $i$. A residue outside its indicated range is
zero. Use the differences $\Delta A_i=A_i-A_{i+1}$,
$\nabla B_i=B_i-sB_{i-1}$. Set
\[
\begin{aligned}
 r_z(X)&=\frac{X-z}{X-sz},& S(u)&=\frac{R(u)}{R(cu)},\\
 M(u,n)&=\frac{(u;p)_n}{(su;p)_n},&
 Q_{m,n}(u)&=\frac{M(u,m+n)}{M(u,m)M(u,n)} .
\end{aligned}
\]
Here $(u;p)_{-n}=1/(p^{-n}u;p)_n$ for $n>0$. We use the transport laws
\begin{equation}\label{eq:qtrans}
 A_{i+1}^\alpha=r_{Z_\alpha}(Z_i)A_i^\alpha,\qquad
 B_{i-1}^\beta=r_{xZ_\beta}(Z_i)B_i^\beta .
\end{equation}
For the first, $P_{i+1}(z)/P_i(z)=(1-z/(sZ_i))/(1-z/Z_i)$ at either
type of step. For the second, $N_i(w)/N_{i-1}(w)$ equals
$R(w/Z_{i-1})$ at a preceding horizontal step and $R(cw/Z_{i-1})$
at a preceding downward step. The increments of $Z_i$ are $x$ and $x/c$,
respectively. Both relations hold at the physical endpoints by clearing
denominators.

Let $G_{11}$ be the following evaluation of $C_1^2$. Expand the slice
differences of each factor. Evaluate the two negative fields on their
horizontal prefixes in increasing slice order, and the two positive fields
on their downward suffixes in decreasing slice order. On ties use center
order and then label order in the chronological word, reversing that order
for the suffix action. The finite Macdonald difference operators retain all
prior shifts of their own alphabet. For each branch with horizontal
multiplicities $m_\beta$ and downward multiplicities $n_\alpha$, insert
\[
                 \prod_{\beta,\alpha}
                 Q_{-m_\beta,n_\alpha}(Z_\alpha/Z_\beta).
\]
This prescription includes all pairs of selected roots. Constants in a
negative or positive field cancel in its own slice difference: their
values are $s^i$ on the prefix and $1$ on the suffix, and the constant
operator shifts no variable. Thus each nonzero branch selects two roots
of each sign.

For $k\le j$, let $E_{NP}$ evaluate the pair $N_j,P_k$ by the indicated
prefix and suffix with its $Q$ factor, including constant choices.
Let $E_{PN}$ instead apply $P_k$ on the prefix and $N_j$ on the suffix,
with their positive and negative shifts. The latter evaluates the
chronological pair. Put $\omega=(1-s)(1-p^{-1})/(1-s/p)$.
\begin{lemma}\label{lem:quadraticinterchange}
For independent boundary roots,
\begin{equation}\label{eq:qmiddle}
 E_{NP}-E_{PN}=
 \omega(1-s)\sum_{k\le h<j}
 \prod_{\substack{k\le r<j\\r\ne h}}R(Z_h/Z_r)\,(\Phi_{h,j}-1),
\end{equation}
where at a horizontal step $\Phi_{h,j}=\prod_{d\ge j,\,D}S(Z_h/Z_d)$,
and at a downward step $\Phi_{h,j}=\prod_{b<j,\,H}S(Z_h/(cZ_b))$.
Here and below $H,D$ restrict a step index to the indicated type.
\end{lemma}
\begin{proof}
Put $f(u)=u/(1-u)$. Direct division gives
$Q_{-1,1}(u)-1=\omega(f(u)-f(su/p))$.
Place all downward roots sufficiently close to zero and all horizontal
roots sufficiently far from zero for one annulus to separate their poles,
also after multiplication by $s$ and $p$. Expand
\[
       \omega\CT_w\big[(P_k(sw)-P_k(pw))N_j(w)\big].
\]
The constant modes cancel. Pairing negative powers of $P_k$ from an inner
downward pole with positive powers of $N_j$ from an outer horizontal
pole gives the product of their signed residues times $Q_{-1,1}-1$.
In the opposite pairing the inner pole of $N_j$ is at $Z_d/p$ and the
outer pole of $P_k$ is at $Z_b$. The geometric sums give
$\omega(f(sZ_d/(pZ_b))-f(Z_d/Z_b))$, with the opposite sign.
The products of the separate means in the two evaluations agree. Thus
the displayed constant term is their difference. Rational continuation
allows arbitrary independent roots, even when a horizontal step in
$[k,j)$ follows a downward step in that interval.

At the cut $j$ the one-body products simplify to
\[
 P_j(sw)N_j(w)=\prod_{d\ge j,\,D}S(w/Z_d),\qquad
 P_j(pw)N_j(w)=\prod_{b<j,\,H}S(w/Z_b).
\]
Changing $j$ to $k$ in $P_j(z)$ multiplies it by
$\prod_{k\le r<j}(1-z/Z_r)/(1-z/(sZ_r))$.
Enlarge the first contour. The value at infinity is $s^{j-k}$, and
the outer residues occur at $Z_h$ for horizontal $h$ in the interval,
with coefficient $(1-s)\prod_{r\ne h}R(Z_h/Z_r)\Phi_{h,j}$.
Shrink the second contour. The value at zero is one and the inner
residues at $Z_h/c$ for downward $h$ have coefficient
$-(1-s)\prod_{r\ne h}R(Z_h/Z_r)\Phi_{h,j}$.
Finally $\sum_h\prod_{r\ne h}R(Z_h/Z_r)=(1-s^{j-k})/(1-s)$,
by partial fractions of $\prod_{k\le r<j}R(w/Z_r)$.
\end{proof}

For centers $i<t$ the outermost fields are the earlier negative and the
later positive, up to commuting operators on a slice. If these choose
$\beta,\alpha$, sum their two slice differences first. Their factor is
\[
   \gamma^2\nabla B_i^\beta\,\Delta A_t^\alpha\,
                      Q_{-1,1}(Z_\alpha/Z_\beta).
\]
For the two middle fields replace $Z_\beta,Z_\alpha$ by
$\widetilde Z_\beta=pZ_\beta,\widetilde Z_\alpha=pZ_\alpha$.
Indeed the outer negative sends $1/Z_\beta$ to $p^{-1}/Z_\beta$;
the outer positive sends $Z_\alpha$ to $pZ_\alpha$.
Both selected roots are present at every middle cut on the appropriate
side. Their joint vacuum factor is the ratio $Q_{-1,1}$ above.
Successive normalized Cauchy factors multiply according to
$M(u,n)M(p^nu,m)=M(u,n+m)$. Applying this identity to each
horizontal--downward pair puts all remaining branch factors at the
two shifted roots, for either assignment of the middle fields.
In particular it also holds when another field selects the same root.
When the middle slices have their negative first, the two evaluations
agree. For the opposite order use Lemma~\ref{lem:quadraticinterchange}
with these shifted roots. Equal centers admit an order with both negative
fields first.

\begin{lemma}\label{lem:quadraticdomino}
Write $h=\beta+1$, $a=\beta+2$, and use the outer shifts at $\beta,a$
in $\Phi_{h,a}$. Then
\begin{multline}\label{eq:qdomino}
 l_1\langle C_2\rangle+
 \frac{l_1^2}{2}\big(\langle C_1^2\rangle-G_{11}\big)\\
 =l_1\gamma\sum_{\beta\,H,\,a\,D}
 B_{\beta+1}^\beta A_a^a\,Q_{-1,1}(Z_a/Z_\beta)\,\Phi_{h,a}.
\end{multline}
\end{lemma}
\begin{proof}
For each middle negative slice $j=t,t-1$ and positive slice $k=i,i+1$
the correction is supported at $k\le h<j$. In particular $\beta<i\le
h<t\le\alpha$. Write
\[
 U_k=\prod_{k\le r<h}R(Z_h/Z_r),\qquad
 V_j=\prod_{h<r<j}R(Z_h/Z_r),\qquad U_{h+1}=V_h=0 .
\]
At fixed $h,\beta,\alpha$ the left position sum is
$\sum_{i=\beta+1}^{h}(U_i-U_{i+1})\nabla B_i^\beta$.
By \eqref{eq:qtrans} the two transport parameters are $a'=Z_h$,
$b'=xZ_\beta$. If $a'\ne b'$ the identity
\begin{equation}\label{eq:qprimitive}
 (1-r_{a'})(1-sr_{b'})
       =\frac{(s-1)a'}{a'-b'}(r_{a'}-r_{b'})
\end{equation}
writes the summand as a forward difference of
$(s-1)a'U_i B_{i-1}^\beta/(a'-b')$, with both boundary values zero.
In the remaining case comparison of the $x$-degrees forces $h=\beta+1$,
and the sum is $B_h^\beta$.

For the right position sum separate $V_j\Phi_{h,j}-V_j$.
The subtracted part uses $B_h^\beta V_t=B_{t-1}^\beta$,
by \eqref{eq:qtrans}, including the extended zero at $t=h$.
Hence after taking the $t,t-1$ difference, multiplying by
$\Delta A_t^\alpha Q_{-1,1}$ and summing, this bare part is
$\langle C_2\rangle/\gamma$: the negative difference for $C_2$
occurs one slice earlier than for $C_1$.
For $W_t=V_t\Phi_{h,t}$ put $W_h=0$. Then
$W_{t-1}=r_{Z_{h+1}}(Z_t)W_t$ up to $t=\alpha$.
In fact, writing $z=Z_h/Z_{t-1}$, the forward quotient $W_t/W_{t-1}$
away from birth is $R(z)$ or $R(cz)$ for $h$ horizontal,
according as $t-1$ is horizontal or downward. For $h$ downward it is
$R(z/c)$ or $R(z)$, respectively. These are the asserted relations
since $Z_t/Z_{t-1}$ is $x$ or $x/c$; the cleared relation at birth has
zero numerator. All crossed steps lie strictly between the shifted
exterior roots. Applying \eqref{eq:qprimitive} a second time leaves only
$\alpha=h+1$ with value $A_\alpha^\alpha\Phi_{h,\alpha}$.
The chronological evaluation uses the negative of \eqref{eq:qmiddle}.
The two orders of distinct centers cancel the factorial $2$, and
$l_1\gamma\omega(1-s)=-1$. This gives \eqref{eq:qdomino}.
\end{proof}

\subsection{The two currents of each sign}
It remains to evaluate $(l_1^2G_{11}+l_2\langle\psi^2C_1\rangle)/2$.
Temporarily give the two selected downward steps distinct labels
$A=\{a,d\}$, their values $Z_\alpha$, and the two horizontal steps
distinct labels $B=\{b,e\}$, their values $xZ_\beta$. Remove the
opposite-sign $Q$ factors, which depend only on these four selections.
Remove also $(l_1\gamma)^2$ and the one-body weights taken at the left
cut for positives and the right cut for negatives. By
\eqref{eq:qtrans} the remaining one-body factors are products of the
$r_z(X)$ over the intervening slice weights $X$. For a set $J$ write
$r_J(X)=\prod_{z\in J}r_z(X)$.

Use $(q_+,V_+,m_+)=(p,s,1)$ and $(q_-,V_-,m_-)=(p^{-1},s^{-1},s)$.
For either sign define
\[
 K_\pm(t)=\frac{(1-t)(1-q_\pm t/V_\pm)}
                    {(1-q_\pm t)(1-t/V_\pm)},\qquad
 E_\pm(u,v)=\frac{V_\pm(u-v)(v-u)}
                  {(V_\pm u-v)(V_\pm v-u)} .
\]
Two distinct selections $u,v$ in chronological order supply
$K_\pm(v/u)$. Indeed dividing the second shifted singleton coefficient by
its unshifted value cancels all factors except that of the first
selected root; reverse the action for a suffix. The coefficient of the
elementary operator $e_2$ divided by the two singletons is $E_\pm$,
by multiplying their two mutual spectator factors and the elementary
prefactor $V_\pm$. Thus the coefficient for $p_2=e_1^2-2e_2$ is
$K_\pm(v/u)+K_\pm(u/v)-2E_\pm(u,v)$. The conversion of a second
power sum to two signed singleton factors contributes
$(V_\pm+1)/(V_\pm-1)$. The product of the two conversion factors times
$-s^2/(s^2-1)^2$ is precisely $\gamma^2$.

Here are the complete densities for one center of slice weight $X$.
For one pair put $d_X(a,b)=(1-r_a(X))(1-sr_b(X))$.
For $u,v$ of one sign write $k=K_\pm(v/u)$, $h=K_\pm(u/v)$,
$r_u=r_u(X)$, $r_v=r_v(X)$, $m=m_\pm$, and set
\begin{align*}
 T_\pm(u,v;X)
 &= (k+h)(1+m^2r_ur_v)-2m
 \begin{cases}r_u h+r_v k&(+),\\r_u k+r_v h&(-),\end{cases}\\
 H_\pm(u,v;X)&=(k+h-2E_\pm(u,v))(1-m^2r_ur_v).
\end{align*}
Indeed both fields at the native slice or both at the displaced slice
retain label order. With just one displaced field, that field is later
for the positive sign and earlier for the negative sign; there are two
labelings. The Adams field makes its slice choice only once.
For two complete pairs at $X$ the density is therefore
\[
 D_X(A,B)=\tfrac12 T_+(A;X)T_-(B;X)
             +\frac{l_2}{2l_1^2}H_+(A;X)H_-(B;X).
\]

\begin{lemma}\label{lem:twoslicetwo}
Through degree two the sum over assignments to centers is invariant
under permutations of the slice weights, after removing the exterior
one-body and opposite-sign factors.
\end{lemma}
\begin{proof}
Consider two consecutive weights $X,Y$. When both pairs are assigned
to these centers, their sum is
\begin{align}
 N(X,Y)={}&r_A(X)D_Y(A,B)+r_B(Y)D_X(A,B)\notag\\
 &+\sum_{\substack{u\in A\\v\in B}}
 K_+(u'/u)K_-(v'/v) r_v(Y)r_{u'}(X)
                   d_X(u,v)d_Y(u',v'),                      \label{eq:qtwoslice}
\end{align}
where primes denote complementary labels. The first term has both pairs
at $Y$, the second both at $X$, and the sum one at each. The factors
joining assignments before or after these two centers depend only on
their union. When one pair occupies the two-center window, its entry is
$r_b(Y)d_X(a,b)+r_a(X)d_Y(a,b)$, symmetric by \eqref{eq:qprimitive}
and continuation.

For two pairs we give the finite calculation of
$N(X,Y)-N(Y,X)$. Clear $\prod_{z\in A\cup B}(X-sz)(Y-sz)$.
The resulting polynomial has degree at most four in each slice
variable: each root occurs at most once per variable in any summand
of \eqref{eq:qtwoslice}. By simultaneous scaling take $A=(1,d)$,
$B=(b,e)$. In the field $\QQ(p,s,d,b,e)$ put
\[
       \frac{l_2}{l_1^2}=\frac{(s+p)(1-p)}{(s-p)(1+p)} .
\]
For $X=-i$, $Y=-j$, the numerators of the difference reduce to the
following entries, in increasing order of $j$,
\[
\begin{array}{c|cccc}
 i\ \backslash\ j&2&3&4&5\\ \hline
 1&0&0&0&0\\
 2&&0&0&0\\
 3&&&0&0\\
 4&&&&0
\end{array}
\]
One can reproduce the table by initializing the two terms on the first
line of \eqref{eq:qtwoslice}, adding its four terms for
$(u,v)=(1,b),(1,e),(d,b),(d,e)$, and subtracting the evaluation
with $i,j$ interchanged. All the denominators at these nodes are
nonzero rational functions. Antisymmetry gives the lower triangle
and the diagonal. Interpolation on five nodes in each variable proves
the identity; thus transposing any two adjacent slice weights preserves
the sum.
\end{proof}

We explain the specialization to repeated selections before permuting
the physical slice weights. If a downward step occurs twice, use the
two positive labels $pZ_\alpha,Z_\alpha$ in that chronological order.
For a repeated horizontal step use $xZ_\beta,pxZ_\beta$, in that
order. In the distinct-label identity divide by the ordered kernel for
the pair and specialize to this string. The kernel has a simple pole,
whereas its reverse and the elementary coefficient $E_\pm$ are
regular. A surviving pair must thus have the indicated order, with
weight one in place of its internal kernel. The divided coefficient
for $p_2$ is also one: a repeated shift occurs once in $e_1^2$ and
never in $e_2$. If both signs repeat, the two normalizing factors
involve independent variables. Consequently the specialization
reproduces the actual Macdonald action. At the second use of one
variable its one-body coefficient is taken at the shifted root,
omitting its own spectator in both coefficients. These weights propagate
as in \eqref{eq:qtrans} at the respective labels on the active cuts.
The common product of opposite-sign factors is legitimate also for
strings, since
\[
     Q_{-m,n}(w)=\prod_{i=0}^{m-1}\prod_{j=0}^{n-1}
                       Q_{-1,1}(p^{j-i}w).
\]

Take a window containing all the physical births $X=xZ_\beta=Z_{\beta+1}$
and deaths $X=Z_\alpha$. An assignment to center $i$ has exterior
transport $\prod_{k<i}r_a(X_k)$ for each positive and
$\prod_{k>i}r_b(X_k)$ for each negative, restricted to this window;
$D_{X_i}$ and $d_{X_i}$ include transport at the center itself.
For a repeated step the shifted copy may be continued beyond the
physical endpoint by its transport law. Every positive string uses
its unshifted copy last and every negative string uses that copy first
in actual slice order. An unshifted copy beyond its allowed range
therefore forces the zero $r_z(z)$, and the continuation introduces
no extra assignment. No physical transport denominator vanishes,
since $Z_i$ has $p$-degree zero and $sz$ has positive $p$-degree.
The specialized symmetry in Lemma~\ref{lem:twoslicetwo} also permits
coincident labels of opposite signs: its only label denominators
apart from transports involve labels of the same sign.

Suppose a selected downward step at $\alpha$ occurs more often than
the horizontal step at $\alpha-1$, counted with multiplicity zero if
absent. Move $X=Z_\alpha$ to the first center. Its unshifted positive
cannot occur later and thus all its copies must be assigned here.
A negative from any other step cannot be used here, since its
unshifted birth now occurs later. This center has insufficient negative
multiplicity. It follows that all nonzero assignments have matching
multiplicities at adjacent horizontal--downward pairs. Indeed an
excess at some downward step exists whenever the multiplicities fail
to match, since both totals are two. In the original slice order,
the first such turn must use all its positive copies there and receives
only its own negatives. The same argument at the next turn pins every
selected field to its corner.

\subsection{Identification of the height weights}
For a consecutive portion of the boundary starting with a horizontal
step $\beta$ and ending with a downward step $a$, let $w_i$ be the
box weights of the marked layers (including $p^{-1}$ for a second
layer) and set $z_i=c/w_i$. Relative to the weight on this portion
alone, the height polarization supplies the factor
\begin{equation}\label{eq:qexternal}
 \prod_i\left(
  \prod_{t<\beta,\,H}\frac{R(z_i/(xZ_t))}{R(z_i/Z_t)}
  \prod_{t>a,\,D}\frac{R(z_i/Z_t)}{R(cz_i/(xZ_t))}
 \right).
\end{equation}
To see this, the boundary term in the height polarization gives
$R(z_i/Z_t)$ at each horizontal--downward turn and its inverse at
each downward--horizontal turn, including the two terminal valleys.
Telescope along the horizontal runs to the left of the portion, and
along the downward runs to its right. The two-layer interaction
depends only on the ratio of the marked weights, so cancels in taking
this relative factor.

At the turn $\beta,a=\beta+1$ selected twice, let $A,B$ be the two
one-body residues there, and $\widetilde A,\widetilde B$ the values for
the second use of the same variable. Both the ordinary square and
the second power sum contribute
$\gamma^2 AB\widetilde A\widetilde B Q_{-2,2}(x)$ before their
respective symmetric-algebra scalars. Their exterior factors are
\eqref{eq:qexternal} for $z_1=Z_a$, $z_2=pZ_a$. On the one-box
diagram their common product reduces to
\[
 \frac{s^2Q_{-2,2}(x)}{R(x)R(px)R(c/x)R(pc/x)}=1,
\]
by $R(z)R(1/(sz))=s$ and the definition of $M$. The coefficient on
this portion is $(l_1^2+l_2)/2=(1-c)(1-pc)/((1-p)(1-p^2))$,
the weight of a height-two box.

For distinct turns write $t=Z_{\alpha_2}/Z_{\alpha_1}$ with the
earlier turn first. The product of their separate one-mark weights
must be multiplied by
\[
 K_+(t)K_-(t)Q_{-1,1}(xt)Q_{-1,1}(x/t)
     =R\big((p^{-1}-1)(1-x)(t+t^{-1})\big).
\]
Here $R$ is extended multiplicatively over a virtual character.
The identity follows by inserting the four fractions and reflecting
the factors in $t^{-1}$ with $R(z)R(1/(sz))=s$. The right side is
precisely the cross term of the height polarization. The two orders of
the marked pairs cancel the factorial in the ordinary square.
The second Adams operation cannot put a single insertion at two
different turns.

Finally the right side of \eqref{eq:qdomino} is supported on the words
$HHD$ and $HDD$. In both cases the two layer parameters in
\eqref{eq:qexternal} are $z_1=Z_{\beta+1}$, $z_2=Z_a$.
For every exterior horizontal $t$ the two residues
$B_{\beta+1}^\beta,A_a^a$ supply $R(Z_\beta/Z_t)/R(Z_a/Z_t)$.
For every exterior downward $t$ they supply $R(Z_a/Z_t)/R(cZ_\beta/Z_t)$.
The additional factor $\Phi_{\beta+1,a}$ multiplies the latter by
$S(Z_{\beta+1}/Z_t)$ on $HHD$, and the former by
$S(Z_{\beta+1}/(cZ_t))$ on $HDD$. Inserting the step increments gives
exactly \eqref{eq:qexternal}. The scalar on each isolated portion can
be seen directly in the following table. In the middle column $z=x$
for the horizontal domino and $z=c/x$ for the vertical one.
\[
\begin{array}{c|c|c}
 \text{portion}&l_1\gamma B_{\beta+1}^\beta A_a^a
       Q_{-1,1}(Z_a/Z_\beta)\Phi_{\beta+1,a}
       &\text{height weight}\\ \hline
 HHD&l_1(1-cz)/(1-pz)&l_1(1-cx)/(1-px)\\
 HDD&l_1(1-cz)/(1-pz)&l_1(1-c^2/x)/(1-pc/x)
\end{array}
\]
For example, set $Z_\beta=1$, $Z_{\beta+1}=x$ and
$Z_a=x^2$ or $x^2/c$, respectively. The horizontal residue at any cut
has own factor $1-s$, factors $R(Z_\beta/Z_t)$ for the other
horizontal steps before that cut and $1/R(cZ_\beta/Z_t)$ for its
downward steps. The downward residue has own factor $s-1$, factors
$1/R(Z_a/Z_t)$ on the horizontal prefix and $R(Z_a/Z_t)$ on the other
downward steps of the suffix. These products give the middle column.
For the last column the two-box row has removable box weight $x^{-1}$,
addable weights $x^{-2},x/c$, and layer sum $1+x^{-1}$; the column
has the same products with $x$ and $c/x$ interchanged, evaluated with
the corresponding two box weights in the given polarization.
Together these calculations account for every size-two reverse plane
partition and prove Theorem~\ref{thm:quadraticdisk}.
